% from ega0-21-fr.tex
\oldpage[0\textsubscript{IV-1}]{179}
  \item[(i)] On sait (21.1.5, (ii)) que
  $\operatorname{rg}_E\Omega^1_{E/M}$ ne change pas quand on remplace $M$ par
  $M(L_0^p)$~; comme
  $L_0^p=K_0^p((T_1^p,\ldots,T_r^p))$ et que
  $[k(K_0^p):k]<+\infty$, on a
  \[
    M(L_0^p)=k(K_0^p)((T_1^p,\ldots,T_r^p))~;
  \]
  de même $\Omega^1_{K_0/k}$ ne change pas quand on y remplace $k$ par
  $k(K_0^p)$, si bien que l'on peut \emph{supposer que $K_0^p\subset k$}, ce
  que nous ferons dans toute la suite. Nous introduirons aussi les corps
  \[
    k'_0=k_0(K_0^p),\qquad
    N'=k'_0((T_1^{p^2},\ldots,T_r^{p^2}))
  \]
  de sorte que l'on a le diagramme de corps
  \[
    \xymatrix{
      K_0\ar[r] & M_0\ar[r] & L_0\ar[r] & E\\
      k\ar[u]\ar[r] & N\ar[u]\ar[r] & M\ar[u]\ar[r] & L\ar[u]\\
      k_0\ar[u]\ar[r] & k'_0\ar[u]\ar[r] & N'\ar[u] &
    }
  \]

  Nous nous proposons d'évaluer la différence $\delta$ du premier et du
  second membre de (21.9.6.2).

  Il résulte de (21.7.3) que l'on a
  \[
    \operatorname{rg}_E\Omega^1_{E/M}
    -\operatorname{rg}_{M_0}\Omega^1_{M_0/M}
    =\deg.\operatorname{tr}_{M_0}E+d(E/M_0,M)
  \]
  et comme $L_0$ est une extension finie de $M_0$, on a
  $\deg.\operatorname{tr}_{L_0}E=\deg.\operatorname{tr}_{M_0}E$. D'autre
  part, une $p$-base de $K_0$ sur $k=k(K_0^p)$ est aussi une $p$-base de
  $M_0$ sur $M$, donc (21.4.5), on a
  \[
    \operatorname{rg}_{M_0}\Omega^1_{M_0/M}
    =\operatorname{rg}_{K_0}\Omega^1_{K_0/k}
  \]
  si bien que l'on a
  \begin{equation}
  \tag{21.9.6.3}
  \label{0.21.9.6.3-fr}
    \delta=d(E/M_0,M)-d(E/K_0,k_0)-r.
  \end{equation}

  Notons maintenant le lemme classique~:

  \begin{lemma}[21.9.6.4]
  \label{0.21.9.6.4-fr}
  Pour tout corps $k$, le corps de séries formelles
  $K=k((T_1,\ldots,T_r))$ est une extension séparable de $k$.
  \end{lemma}

  \begin{proof}
  Rappelons rapidement la démonstration de ce lemme pour être complet. Il
  suffit (Bourbaki, \emph{Alg.}, chap.~VIII, \S~7,
  n\textsuperscript{o}~3, démonstration du th.~1) de prouver que pour toute
  extension finie $k'$ de $k$, $K\otimes_k k'$ est sans élément nilpotent~;
  mais si l'on pose $A=k[[T_1,\ldots,T_r]]$ et $S=A-\{0\}$,
  $K\otimes_k k'$ est égal à $S^{-1}(A\otimes_k k')$, et
  $A\otimes_k k'$ s'identifie canoniquement à l'anneau intègre
  $A'=k'[[T_1,\ldots,T_r]]$ et $A$ à un sous-anneau de $A'$, d'où la
  conclusion.
  \end{proof}

  Utilisant ce lemme et (21.6.2, (ii)), on a $d(L_0/K_0,k_0)=0$, et la
  formule (21.6.4.4) donne donc
  \[
    d(E/K_0,k_0)=d(E/L_0,k_0).
  \]

  